\documentclass{article}
\usepackage{amsmath,amsthm,amssymb}

\newtheorem{theorem}{Theorem}

\title{A Minimal Fixed-Point Note}
\author{Synthetic evaluation fixture}
\date{}

\begin{document}
\maketitle

\section{Contractions and completeness}

\begin{theorem}[Contraction mapping theorem]\label{th:contraction}
Let $(X,d)$ be a nonempty complete metric space. Suppose that $T:X\to X$ satisfies
\[
  d(Tx,Ty)\leq q\,d(x,y)
\]
for all $x,y\in X$ and some constant $0\leq q<1$. Then $T$ has a unique fixed point $x_*$, and for every $x_0\in X$ the sequence defined by $x_{n+1}=T(x_n)$ converges to $x_*$.
\end{theorem}

\begin{proof}
Starting from $x_0\in X$, define $x_{n+1}=T(x_n)$. Iterating the contraction estimate gives
\[
  d(x_{n+1},x_n)\leq q^n d(x_1,x_0).
\]
Hence, for $m>n$, the triangle inequality and a geometric-series estimate give
\[
  d(x_m,x_n)
  \leq \sum_{k=n}^{m-1}d(x_{k+1},x_k)
  \leq \frac{q^n}{1-q}d(x_1,x_0).
\]
Thus $(x_n)$ is Cauchy. Completeness supplies a limit $x_*\in X$. Since the contraction estimate implies continuity of $T$, passing to the limit in $x_{n+1}=T(x_n)$ gives $T(x_*)=x_*$. If $y_*$ is another fixed point, then
\[
  d(x_*,y_*)=d(Tx_*,Ty_*)\leq qd(x_*,y_*),
\]
and $q<1$ forces $d(x_*,y_*)=0$.
\end{proof}

\end{document}
